\documentclass[11pt,a4paper]{article}

\usepackage[margin=1in]{geometry}
\usepackage{amsmath,amssymb,amsthm,bm,graphicx}
\usepackage{enumitem}
\usepackage{hyperref}
\usepackage{url}
\makeatletter
\g@addto@macro{\UrlBreaks}{\do\.\do\_\do\-}
\makeatother
\newcommand{\lmod}[1]{\href{http://www.lemma.cn/lean/?module=#1}{\nolinkurl{#1}}}

\hypersetup{
  colorlinks=true,
  linkcolor=blue,
  citecolor=blue,
  urlcolor=blue
}

\newtheorem{theorem}{Theorem}[section]
\newtheorem{lemma}[theorem]{Lemma}
\newtheorem{proposition}[theorem]{Proposition}
\newtheorem{definition}[theorem]{Definition}
\newtheorem{remark}[theorem]{Remark}

\newcommand{\R}{\mathbb{R}}
\newcommand{\N}{\mathbb{N}}
\newcommand{\Tensor}{\operatorname{Tensor}}
\newcommand{\rot}{\mathrm{R}}
\newcommand{\eye}{\mathrm{I}}
\newcommand{\softmax}{\mathrm{softmax}}
\newcommand{\concat}{\mathbin{+\mkern-4mu+}}
\newcommand{\Rel}{\mathrm{Rel}}
\emergencystretch=3em

\title{Lean~4-Checked Identities for\\One-Dimensional Rotary Position Embeddings}
\author{LiZhi Zhou\\[0.3em]
\small Project artifact: \href{https://github.com/math-proof/lemma}{github.com/math-proof/lemma}\\[0.4em]
\small
Relative:
\href{http://www.lemma.cn/lean/?module=Tensor.DotSoftmaxDivDot_Stack_TDot.eq.Stack_Div_SumExp.of.Eq_Stack_Mul}{softmax},
\quad
pairing:
\href{http://www.lemma.cn/lean/?module=Tensor.DotDotSRotaryMatrix.eq.Dot_DotRotaryMatrixSub}{$q^\top R_{n-m}k$},
\quad
layouts:
\href{http://www.lemma.cn/lean/?module=Tensor.RotaryMatrix'.eq.DotDot_RotaryMatrix}{$R'=P^\top RP$}\\[0.25em]
realization:
\href{http://www.lemma.cn/lean/?module=Tensor.DotRotaryMatrix.eq.AddMulS}{kernel},
\quad
additive:
\href{http://www.lemma.cn/lean/?module=Tensor.DotRotaryMatrixS.eq.RotaryMatrixAdd}{group law},
\quad
orthogonal:
\href{http://www.lemma.cn/lean/?module=Tensor.DotT_RotaryMatrix.eq.Eye}{$R^\top R=I$}}
\date{}

\begin{document}
\maketitle

\begin{abstract}
Rotary position embedding (RoPE) encodes a token index by rotating query and key vectors in two-dimensional planes.
The identities that make this useful---a kernel evaluation formula, orthogonality, angle addition, and the reduction of attention to a relative offset---are standard, but they are usually written for $2\times 2$ blocks or for a fixed geometric frequency schedule, and they are not machine-checked.
We develop one-dimensional RoPE in the \texttt{lemma} library in Lean~4.
Two matrices \(\rot(\alpha),\rot'(\alpha)\in\R^{(d+d)\times(d+d)}\) are defined for an arbitrary angle vector \(\alpha\in\R^d\) (here \(d\) is already the half-dimension).
The unprimed matrix is the even/odd split-half layout of Hugging Face Transformers~\cite{huggingface_rope_utils}; the primed matrix is Su et al.'s original interleaved block-diagonal layout~\cite{su2021roformer}.
They are conjugate by a fixed even/odd gather \(\boldsymbol{P}\):
\(\rot'(\alpha)=\boldsymbol{P}^\top\rot(\alpha)\boldsymbol{P}\),
and entrywise
\(\rot'(\alpha)_{i,j}=\rot(\alpha)_{i^{\mathrm{toSplit}},j^{\mathrm{toSplit}}}\).
We prove, with no \texttt{sorry}:
the realization
\(\rot(\alpha)\,x = x\odot\cos(\alpha\concat\alpha)+((-x_1)\concat x_0)\odot\sin(\alpha\concat\alpha)\);
the orthogonal identities \(\rot(\alpha)^\top\rot(\alpha)=\eye_{d+d}\) and \(\rot(-\alpha)=\rot(\alpha)^\top\);
the free-angle group law \(\rot(\alpha)\rot(\beta)=\rot(\alpha+\beta)\);
the RoFormer pairing
\((\rot(\alpha)\,q)^\top(\rot(\beta)\,k)=q^\top\rot(\beta-\alpha)\,k\);
the same transpose product on the primed layout,
\(\rot'(\alpha)^\top\rot'(\beta)=\rot'(\beta-\alpha)\);
and, under the linear hypothesis \(\theta_i=i\,\tau\), the relative form of softmax attention, in which each unnormalized weight depends on \(\rot(\theta_k-\theta_i)\) rewritten as \(\rot(\theta_{k-i})\) or \(\rot(\theta_{i-k})^\top\).
Interactive statements are public on lemma.cn~\cite{lemma_rope_rel,lemma_rope_pair,lemma_rope_conj,lemma_rope_real,lemma_rope_add,lemma_rope_orth}.
\end{abstract}

\section{Introduction}
\label{sec:intro}

Transformers need a representation of order~\cite{vaswani2017attention}.
Rotary position embedding (RoPE)~\cite{su2021roformer,su2024roformer} supplies one by rotating each pair of channels of the query and key through an angle that depends on the token index.
The construction is now the default positional bias in open language models~\cite{touvron2023llama}.
Its algebraic content is classical: plane rotations form an abelian subgroup of \(\mathrm{SO}(2)\), the transpose is the inverse, and products add angles.
RoFormer records the consequence that matters for attention,
\begin{equation}
  \label{eq:roformer-rel}
  \bigl(\rot_{\Theta,m}\,q\bigr)^\top\bigl(\rot_{\Theta,n}\,k\bigr)
  = q^\top \rot_{\Theta,n-m}\,k,
\end{equation}
and introduces the relative matrix as a transpose product
\(\rot_{\Theta,n-m}=(\rot_{\Theta,m})^\top\rot_{\Theta,n}\).
Equation~\eqref{eq:roformer-rel} is machine-checked in Lean~4 as
\lmod{Tensor.DotDotSRotaryMatrix.eq.Dot_DotRotaryMatrixSub};
the proof is recorded in Section~\ref{sec:pair}.
Later expositions isolate the two-by-two rules \(\rho(\theta)^\top=\rho(-\theta)\) and \(\rho(\theta)\rho(\theta')=\rho(\theta+\theta')\)~\cite{barbero2025round}.

Those identities are easy to check by hand on a single plane and then lift through a block-diagonal direct sum.
They are less pleasant to check on the layout that implementations actually multiply: the embedding is split into a first half and a second half, and the rotation of plane \(p\) mixes coordinates \(p\) and \(p+d\) rather than \(2p\) and \(2p+1\).
The efficient evaluation
\begin{equation}
  \label{eq:impl-intro}
  \rot(\alpha)\,x
  = x \odot \cos(\alpha\concat\alpha)
  + \mathrm{rotateHalf}(x)\odot\sin(\alpha\concat\alpha)
\end{equation}
is the form that appears in Hugging Face kernels~\cite{huggingface_rope_utils,touvron2023llama}; it is not the interleaved matrix in which~\eqref{eq:roformer-rel} is written.
A gap between the RoFormer matrix and the kernel formula is a natural place for a machine-checked development.

This paper records such a development in Lean~4~\cite{moura2021lean4}, on top of the \texttt{lemma} tensor library~\cite{lemma_repo} in which shapes are lists of natural numbers and matrix multiplication is a typed \texttt{Dot} instance.
The objects we study are two functions
\[
  \rot,\;\rot' : \R^d \to \R^{(d+d)\times(d+d)}
\]
of a free angle vector:
\texttt{rotaryMatrix} is the Hugging Face split-half matrix, and
\texttt{rotaryMatrix'} is Su et al.'s original interleaved matrix.
No frequency schedule is built into \(\rot\).
The geometric RoPE angles \(\theta_{i,p}=\lambda\,i/b^{p/d}\) appear only later, and only as the special case \(\theta_i=i\,\tau\) of a constant frequency vector \(\tau\in\R^d\).
That hypothesis is the arithmetic needed for relative indexing; YaRN-style scalings and other learned or interpolated frequencies~\cite{peng2023yarn} are the same lemma with a different \(\tau\).

\paragraph{Contributions.}
We machine-check, with no \texttt{sorry}, the following groups of one-dimensional RoPE facts.
Each public name below is a module path in the library; the path \emph{is} the statement.

\begin{enumerate}[leftmargin=1.6em,itemsep=0.4em]
  \item \textbf{Two layouts.}
    \lmod{Tensor.RotaryMatrix.eq.AppendHstackSMulSEye}~\cite{lemma_rope_hf_def}
    is Hugging Face \texttt{rotate\_half}~\cite{huggingface_rope_utils};
    \lmod{Tensor.RotaryMatrix'.eq.Stack_Ite_IteS}~\cite{lemma_rope_su_def}
    is the RoFormer matrix~\cite{su2021roformer}.
    They are conjugate by the even/odd gather \(\boldsymbol{P}=\mathrm{interleave}\,d\):
    \lmod{Tensor.RotaryMatrix'.eq.DotDot_RotaryMatrix}~\cite{lemma_rope_conj}
    states \(\rot'(\alpha)=\boldsymbol{P}^\top\rot(\alpha)\boldsymbol{P}\),
    and
    \lmod{Tensor.GetRotaryMatrix'.eq.GetRotaryMatrix}~\cite{lemma_rope_entry}
    states the entrywise form.
  \item \textbf{Realization.}
    \lmod{Tensor.DotRotaryMatrix.eq.AddMulS}~\cite{lemma_rope_real}:
    equation~\eqref{eq:impl-intro} for the Hugging Face matrix.
  \item \textbf{Orthogonality.}
    \lmod{Tensor.DotT_RotaryMatrix.eq.Eye}~\cite{lemma_rope_orth}:
    \(\rot(\alpha)^\top\rot(\alpha)=\eye_{d+d}\),
    together with \(\rot(-\alpha)=\rot(\alpha)^\top\) and \(\rot(-\alpha)\rot(\alpha)=\eye_{d+d}\).
  \item \textbf{Additivity.}
    \lmod{Tensor.DotRotaryMatrixS.eq.RotaryMatrixAdd}~\cite{lemma_rope_add}:
    \(\rot(\alpha)\rot(\beta)=\rot(\alpha+\beta)\)
    for arbitrary angle vectors, not only for integer token indices.
  \item \textbf{Pairing.}
    \lmod{Tensor.DotDotSRotaryMatrix.eq.Dot_DotRotaryMatrixSub}~\cite{lemma_rope_pair}:
    equation~\eqref{eq:roformer-rel} for free angle vectors,
    \((\rot(\alpha)\,q)^\top(\rot(\beta)\,k)=q^\top\rot(\beta-\alpha)\,k\).
    The same transpose product on the RoFormer layout is
    \lmod{Tensor.DotT.eq.RotaryMatrix'Sub}~\cite{lemma_rope_su_sub}:
    \(\rot'(\alpha)^\top\rot'(\beta)=\rot'(\beta-\alpha)\).
  \item \textbf{Relative attention.}
    \lmod{Tensor.DotSoftmaxDivDot_Stack_TDot.eq.Stack_Div_SumExp.of.Eq_Stack_Mul}~\cite{lemma_rope_rel}:
    after applying \(\rot(\theta_i)\) to every query and key,
    each unnormalized attention weight depends on \(\rot(\theta_k-\theta_i)\),
    and under \(\theta_i=i\,\tau\) this matrix is
    \(\rot(\theta_{k-i})\) or \(\rot(\theta_{i-k})^\top\).
\end{enumerate}

The mathematics is not new.
What is packaged here is a pair of formal objects---the RoFormer matrix and the Hugging Face matrix---a checked conjugation between them, a frequency hypothesis no stronger than linearity in the index, and a complete chain from the Hugging Face definition to the softmax identity.

\section{Related work}
\label{sec:related}

RoFormer~\cite{su2021roformer,su2024roformer} introduces RoPE, writes the block-diagonal matrix of two-by-two rotations on interleaved pairs \((2p,\,2p+1)\), and derives~\eqref{eq:roformer-rel}.
The paper emphasizes orthogonality and the relative-offset form of the inner product.
It introduces the relative matrix as a transpose product
\(\rot_{\Theta,n-m}=(\rot_{\Theta,m})^\top\rot_{\Theta,n}\),
rather than as the free-angle composition \(\rot(\alpha)\rot(\beta)=\rot(\alpha+\beta)\).
Hugging Face Transformers uses the conjugate split-half layout and evaluates it by \texttt{rotate\_half}~\cite{huggingface_rope_utils}.
The two published matrices are~\eqref{eq:rot-su} and~\eqref{eq:rot-hf} below; this paper machine-checks both and the conjugation.

That additive law is already stated in later theoretical work.
Liu et al.~\cite{liu2025rethinking} call it \emph{exponential additivity}:
if \(XY=YX\), then
\(\exp(X)\exp(Y)=\exp(X+Y)\).
With the Lie parameterization \(\rot_{\boldsymbol{x}}=\exp(\sum_i x_i B_i)\) and commuting generators, this is exactly
\begin{equation}
  \label{eq:liu-add}
  \rot_{\boldsymbol{u}}\rot_{\boldsymbol{v}} = \rot_{\boldsymbol{u}+\boldsymbol{v}},
\end{equation}
and they record the companion identity
\(\rot_{\boldsymbol{u}}^\top\rot_{\boldsymbol{v}}=\rot_{\boldsymbol{v}-\boldsymbol{u}}\).
Zhang et al.~\cite{zhang2025grape} write the same fact as the one-parameter subgroup law
\(\mathbf{G}(n+m)=\mathbf{G}(n)\mathbf{G}(m)\)
and recover RoPE as multiplicative GRAPE.
Barbero et al.~\cite{barbero2025round} use the integer special case
\(\rho(g)^i=\rho(ig)\), so
\(\rot^i\rot^j=\rot^{i+j}\).
Context-extension methods such as YaRN~\cite{peng2023yarn} keep this group law and change only the frequencies.
The present paper does not claim~\eqref{eq:liu-add} as new mathematics;
it machine-checks the identity for the Hugging Face matrix and an arbitrary pair of angle vectors, and transports the transpose product to the RoFormer matrix.
Implementation notes~\cite{touvron2023llama,huggingface_rope_utils} evaluate RoPE through~\eqref{eq:impl-intro} rather than a dense \((d+d)\times(d+d)\) multiply.

Relative attention more broadly includes Shaw et al.~\cite{shaw2018relative} and ALiBi~\cite{press2022alibi}.
We do not compare methods; we only certify the algebra of the rotary special case.

On the formal side, mathlib~\cite{mathlib2020,mathlib} contains the trigonometric addition formulae that sit under our real and vector lemmas, but not a RoPE matrix or an attention identity.

\section{Tensors and notation}
\label{sec:prelim}

The ambient type is a shaped tensor \(\Tensor\,\R\,s\) with \(s:\mathrm{List}\,\N\).
A vector of length \(d\) is a tensor of shape \([d]\); a matrix of size \((d+d)\times(d+d)\) has shape \([d+d,\,d+d]\).
Addition, multiplication, cosine and sine are pointwise.
Matrix multiplication is written \(A @ B\) and is an instance of a \texttt{Dot} class whose result shape is \(\mathrm{matmul\_shape}\).
In particular, the inner product of two \(1\)-tensors is again written \(@\); the library wraps that scalar as \(\Tensor\,\R\,[]\).
Orthogonality is the equation \(A^\top @ A = \eye\).

Throughout, Lean's parameter \(d\) is the \emph{half}-dimension: the embedding width is \(d+d\).
This matches the even/odd (split-half) pairing and avoids a side condition that the width is even.
We write \(\concat\) for concatenation of \(1\)-tensors and \(\odot\) for pointwise product.
Index notation \(\theta[i]\) is \texttt{GetElem} along the leading axis.
For \(k,t:\mathrm{Fin}\,n\) with \(k\ge t\), the difference \(k-t\) agrees with subtraction of values (\lmod{Nat.ValSub.eq.SubValS.of.Ge}).

We write \(\rot(\alpha)\) for the Lean function \texttt{rotaryMatrix} (Hugging Face) and \(\rot'(\alpha)\) for \texttt{rotaryMatrix'} (RoFormer).
The even/odd gather is \(\boldsymbol{P}=\mathrm{interleave}\,d\).

\section{The rotary matrix}
\label{sec:def}

The library defines two rotary matrices on the same angle vector \(\alpha\in\R^d\).
They differ only by which coordinates form a plane.
The displayed arrays below are the LaTeX from the corresponding lemma docstrings; both write \({\boldsymbol{R}}^{d}_{\Theta,m}\) for the published matrix, while Lean's \(d\) is the half-dimension, so \(\alpha\) plays the role of \((m\theta_1,\ldots,m\theta_{d/2})\) and the matrix has shape \([d+d,\,d+d]\).

\subsection{Hugging Face split-half matrix}
\label{sec:def-hf}

\begin{definition}[Split-half rotary matrix]
\label{def:rot}
Let \(\alpha\in\R^d\) and let \(\eye_d\) be the \(d\times d\) identity.
Define
\begin{equation}
  \label{eq:rot-def}
  \rot(\alpha)
  \;=\;
  \begin{pmatrix}
    \eye_d \odot \cos\alpha & -(\eye_d \odot \sin\alpha) \\
    \eye_d \odot \sin\alpha & \phantom{-}\eye_d \odot \cos\alpha
  \end{pmatrix}
  \in \R^{(d+d)\times(d+d)},
\end{equation}
where \(\odot\) broadcasts the vector \(\cos\alpha\) (resp.\ \(\sin\alpha\)) onto the diagonal of each block.
In the library this is a vertical concatenation of two horizontal stacks
(\lmod{Tensor.RotaryMatrix.eq.AppendHstackSMulSEye}~\cite{lemma_rope_hf_def}).
\end{definition}

This is the even/odd layout of Hugging Face Transformers~\cite{huggingface_rope_utils}: plane \(p\) is the pair \((p,\,p+d)\), matching \texttt{rotate\_half}.
The lemma docstring records the same matrix in the published expanded form
\begin{equation}
\label{eq:rot-hf}
\resizebox{\linewidth}{!}{\(\displaystyle
{\boldsymbol{R}}^{d}_{\Theta,m}=\begin{pmatrix}
\cos m\theta_{1} & 0 & \cdots & 0 & -\sin m\theta_{1} & 0 & \cdots & 0 \\
0 & \cos m\theta_{2} & \cdots & 0 & 0 & -\sin m\theta_{2} & \cdots & 0 \\
\vdots & \vdots & \ddots & \vdots & \vdots & \vdots & \ddots & \vdots \\
0 & 0 & \cdots & \cos m\theta_{d/2} & 0 & 0 & \cdots & -\sin m\theta_{d/2} \\
\sin m\theta_{1} & 0 & \cdots & 0 & \cos m\theta_{1} & 0 & \cdots & 0 \\
0 & \sin m\theta_{2} & \cdots & 0 & 0 & \cos m\theta_{2} & \cdots & 0 \\
\vdots & \vdots & \ddots & \vdots & \vdots & \vdots & \ddots & \vdots \\
0 & 0 & \cdots & \sin m\theta_{d/2} & 0 & 0 & \cdots & \cos m\theta_{d/2}
\end{pmatrix}
\)}
\end{equation}
Equation~\eqref{eq:rot-def} is~\eqref{eq:rot-hf} in block form.

Acting on a split vector \(x=x_0\concat x_1\) with \(x_0,x_1\in\R^d\), the definition is the simultaneous plane rotation
\begin{equation}
  \label{eq:plane}
  \begin{pmatrix} y_{0,p} \\ y_{1,p} \end{pmatrix}
  =
  \begin{pmatrix}
    \cos\alpha_p & -\sin\alpha_p \\
    \sin\alpha_p & \phantom{-}\cos\alpha_p
  \end{pmatrix}
  \begin{pmatrix} x_{0,p} \\ x_{1,p} \end{pmatrix}.
\end{equation}

\subsection{RoFormer interleaved matrix}
\label{sec:def-su}

\begin{definition}[Interleaved rotary matrix]
\label{def:rotp}
\lmod{Tensor.RotaryMatrix'.eq.Stack_Ite_IteS}~\cite{lemma_rope_su_def}.
Su et al.'s original rotary matrix~\cite{su2021roformer} acts on interleaved pairs \((2p,\,2p+1)\) (one-based \((2i-1,\,2i)\)).
The lemma docstring records
\begin{equation}
\label{eq:rot-su}
\resizebox{\linewidth}{!}{\(\displaystyle
{\boldsymbol{R}}^{d}_{\Theta,m}=\begin{pmatrix}
\cos m\theta_{1} & -\sin m\theta_{1} & 0 & 0 & \cdots & 0 & 0 \\
\sin m\theta_{1} & \cos m\theta_{1} & 0 & 0 & \cdots & 0 & 0 \\
0 & 0 & \cos m\theta_{2} & -\sin m\theta_{2} & \cdots & 0 & 0 \\
0 & 0 & \sin m\theta_{2} & \cos m\theta_{2} & \cdots & 0 & 0 \\
\vdots & \vdots & \vdots & \vdots & \ddots & \vdots & \vdots \\
0 & 0 & 0 & 0 & \cdots & \cos m\theta_{d/2} & -\sin m\theta_{d/2} \\
0 & 0 & 0 & 0 & \cdots & \sin m\theta_{d/2} & \cos m\theta_{d/2}
\end{pmatrix}
\)}
\end{equation}
We write \(\rot'(\alpha)\) for this matrix.
On an even row \(i=2p\) the nonzero entries are \(\cos\alpha_p\) on the diagonal and \(-\sin\alpha_p\) immediately to the right; on the next row they are \(\sin\alpha_p\) immediately to the left and \(\cos\alpha_p\) on the diagonal
(\lmod{Tensor.GetRotaryMatrix'.eq.Ite_IteS}).
\end{definition}

\subsection{Comparison}
\label{sec:compare}

The two layouts differ by a fixed even/odd gather, independent of \(\alpha\).

\begin{definition}[Even/odd gather]
\label{def:P}
\lmod{Tensor.Interleave.eq.AppendStackS_Delta}.
Let \(\boldsymbol{P}:=\mathrm{interleave}\,d\in\R^{(d+d)\times(d+d)}\) be the permutation matrix from the lemma docstring
\begin{equation}
\label{eq:P}
{\boldsymbol{P}}=\begin{pmatrix}
1 & 0 & 0 & 0 & \cdots & 0 & 0 \\
0 & 0 & 1 & 0 & \cdots & 0 & 0 \\
\vdots & \vdots & \vdots & \vdots & \ddots & \vdots & \vdots \\
0 & 0 & 0 & 0 & \cdots & 1 & 0 \\
0 & 1 & 0 & 0 & \cdots & 0 & 0 \\
0 & 0 & 0 & 1 & \cdots & 0 & 0 \\
\vdots & \vdots & \vdots & \vdots & \ddots & \vdots & \vdots \\
0 & 0 & 0 & 0 & \cdots & 0 & 1
\end{pmatrix}
\end{equation}
so \((\boldsymbol{P}x)_i=x_{2i}\) and \((\boldsymbol{P}x)_{i+d}=x_{2i+1}\).
\end{definition}

The matrix \(\boldsymbol{P}\) is orthogonal:
\(\boldsymbol{P}^\top@\boldsymbol{P}=\eye_{d+d}\)
(\lmod{Tensor.DotTInterleave.eq.Eye})
and
\(\boldsymbol{P}@\boldsymbol{P}^\top=\eye_{d+d}\)
(\lmod{Tensor.Dot_TInterleave.eq.Eye}).

\begin{theorem}[Conjugation]
\label{thm:conj}
\lmod{Tensor.RotaryMatrix'.eq.DotDot_RotaryMatrix}~\cite{lemma_rope_conj}.
For every \(\alpha\in\R^d\),
\begin{equation}
  \label{eq:conj}
  \rot'(\alpha)
  \;=\;
  \bigl((\mathrm{interleave}\,d)^\top @ \rot(\alpha)\bigr) @ \mathrm{interleave}\,d.
\end{equation}
\end{theorem}

Thus \(\rot'(\alpha)=\boldsymbol{P}^\top\rot(\alpha)\boldsymbol{P}\): an interleaved vector is gathered to the split-half layout, rotated by the Hugging Face matrix, and scattered back.
The four quadrants of~\eqref{eq:rot-hf} become the \(2\times 2\) blocks of~\eqref{eq:rot-su}.

Entrywise, the same fact is a reindexing of coordinates.

\begin{lemma}
\label{lem:toSplit}
\lmod{Fin.ToSplit.eq.Ite_Div_2}.
For \(i:\mathrm{Fin}\,(d+d)\),
\begin{equation}
  i^{\mathrm{toSplit}}
  \;=\;
  \begin{cases}
    i/2 & \text{if \(i\) is even},\\
    i/2+d & \text{if \(i\) is odd}.
  \end{cases}
\end{equation}
Even interleaved indices land in the first half; odd ones land in the second half.
\end{lemma}

\begin{theorem}[Entrywise comparison]
\label{thm:entry}
\lmod{Tensor.GetRotaryMatrix'.eq.GetRotaryMatrix}~\cite{lemma_rope_entry}.
For every \(\alpha\in\R^d\) and \(i,j:\mathrm{Fin}\,(d+d)\),
\begin{equation}
  \label{eq:entry}
  \rot'(\alpha)_{i,j}
  \;=\;
  \rot(\alpha)_{i^{\mathrm{toSplit}},\,j^{\mathrm{toSplit}}}.
\end{equation}
\end{theorem}

The four cases of~\eqref{eq:entry} are the four quadrant lemmas
\lmod{Tensor.GetRotaryMatrix.eq.MulCos_Delta.of.Lt.Lt},
\lmod{Tensor.GetRotaryMatrix.eq.MulCos_Delta.of.Ge.Ge},
\lmod{Tensor.GetRotaryMatrix.eq.MulSin_Delta.of.Ge.Lt},
and
\lmod{Tensor.GetRotaryMatrix.eq.MulNegSin_Delta.of.Lt.Ge}.

Kernel realization (Section~\ref{sec:real}) is a statement about \(\rot\), because that is the layout \texttt{rotate\_half} implements.
The RoFormer pairing chain in Section~\ref{sec:pair} is written for \(\rot\); the matching matrix identity for \(\rot'\) is Theorem~\ref{thm:dotTp}.

\section{Realization}
\label{sec:real}

A dense multiply by~\eqref{eq:rot-def} is not how Hugging Face RoPE is run.
The kernel form uses one cosine, one sine, and a half-swap.

\begin{theorem}[Realization]
\label{thm:real}
\lmod{Tensor.DotRotaryMatrix.eq.AddMulS}~\cite{lemma_rope_real}.
Let \(\alpha\in\R^d\) and \(x\in\R^{d+d}\), and write
\[
  x_0 := \mathrm{cast}(x_{[:d]}),\qquad
  x_1 := \mathrm{cast}(x_{[d:]}),
\]
where the casts transport the dependent slice shapes along
\lmod{List.ConsLengthSlice.eq.List}.
Then
\begin{equation}
  \label{eq:real}
  \rot(\alpha)\,@\,x
  \;=\;
  x \odot \bigl(\cos(\alpha\concat\alpha)\bigr)
  \;+\;
  \bigl((-x_1)\concat x_0\bigr) \odot \bigl(\sin(\alpha\concat\alpha)\bigr).
\end{equation}
\end{theorem}

The proof first identifies \(x=x_0\concat x_1\) up to the library's shape-cast equality \(\simeq\), then applies the already-split identity
\lmod{Tensor.DotRotaryMatrix.eq.AddMulSAppend}:
unfold Definition~\ref{def:rot}, multiply the four blocks against \(x_0\concat x_1\), and reassemble the concatenations.
The identity \((-x_1)\concat x_0\) is exactly \texttt{rotate\_half} in Hugging Face Transformers~\cite{huggingface_rope_utils}.
Theorem~\ref{thm:real} is therefore the certificate that the kernel and the Hugging Face matrix are the same map.
The interleaved matrix \(\rot'\) is recovered afterwards by Theorem~\ref{thm:conj}.

\section{Orthogonality}
\label{sec:orth}

Each pair of coordinates \((x_{0,p},x_{1,p})\) is rotated by the standard \(2\times 2\) matrix
\begin{equation}
  \label{eq:so2}
  \rho(\alpha_p)
  \;=\;
  \begin{pmatrix}
    \cos\alpha_p & -\sin\alpha_p \\
    \sin\alpha_p & \phantom{-}\cos\alpha_p
  \end{pmatrix},
\end{equation}
which is an element of \(\mathrm{SO}(2)\): the group of orientation-preserving linear maps of the plane that keep Euclidean length, equivalently
\(\{A\in\R^{2\times 2}:A^\top A=\eye_2,\;\det A=1\}\).
The split-half matrix \(\rot(\alpha)\) applies one such \(\rho(\alpha_p)\) independently in every plane \(p\), so \(\rot(\alpha)\) itself lies in \(\mathrm{SO}(d+d)\).

\begin{lemma}
\label{lem:neg-T}
\lmod{Tensor.RotaryMatrixNeg.eq.TRotaryMatrix}.
For every \(\alpha\in\R^d\),
\begin{equation}
  \rot(-\alpha) = \rot(\alpha)^\top.
\end{equation}
\end{lemma}

The proof is the even/odd identities \(\cos(-\alpha)=\cos\alpha\) and \(\sin(-\alpha)=-\sin\alpha\)
(\lmod{Tensor.CosNeg.eq.Cos}, \lmod{Tensor.SinNeg.eq.NegSin}),
together with the block-transpose law for a \(2\times 2\) stack of horizontal concatenations
(\lmod{Tensor.TAppendHstackS.eq.AppendHstackSTS}).

\begin{lemma}
\label{lem:R0}
\lmod{Tensor.RotaryMatrix0.eq.Eye}.
\(\rot(0)=\eye_{d+d}\).
The RoFormer layout satisfies the same identity:
\lmod{Tensor.RotaryMatrix'0.eq.Eye} states \(\rot'(0)=\eye_{d+d}\).
\end{lemma}

This is \(\cos 0=1\), \(\sin 0=0\)
(\lmod{Tensor.Cos0.eq.One}, \lmod{Tensor.Sin0.eq.Zero})
and the block identity
\lmod{Tensor.AppendHstackS.eq.Eye}.

\begin{theorem}[Orthogonality]
\label{thm:orth}
\lmod{Tensor.DotT_RotaryMatrix.eq.Eye}~\cite{lemma_rope_orth}.
For every \(\alpha\in\R^d\),
\begin{equation}
  \rot(\alpha)^\top @ \rot(\alpha) = \eye_{d+d}.
\end{equation}
\end{theorem}

\begin{proof}
The difference identity of Section~\ref{sec:add} gives
\(\rot(\alpha)^\top @ \rot(\alpha)=\rot(\alpha-\alpha)=\rot(0)\).
Lemma~\ref{lem:R0} finishes the argument.
\end{proof}

A corollary, combining Lemma~\ref{lem:neg-T} and Theorem~\ref{thm:orth}, is the left inverse in the opposite order
(\lmod{Tensor.DotRotaryMatrixNeg.eq.Eye}):
\begin{equation}
  \rot(-\alpha)\,@\,\rot(\alpha) = \eye_{d+d}.
\end{equation}

\section{Additivity}
\label{sec:add}

\begin{theorem}[Transpose product]
\label{thm:dotT}
\lmod{Tensor.DotT.eq.RotaryMatrixSub}.
For all \(\alpha,\beta\in\R^d\),
\begin{equation}
  \rot(\alpha)^\top @ \rot(\beta) = \rot(\beta-\alpha).
\end{equation}
\end{theorem}

This is the free-angle form of RoFormer's definition of \(\rot_{\Theta,n-m}\), stated on the Hugging Face matrix.
The proof expands both sides of~\eqref{eq:rot-def}, multiplies the four-by-four pattern of blocks
(\lmod{Tensor.DotAppendSHstackS.eq.AppendHstackSAddSDotS},
\lmod{Tensor.DotMulSEye.eq.MulEye}),
and applies the tensor cosine- and sine-subtraction formulae
\begin{align}
  \cos(\beta-\alpha)&=\cos\alpha\odot\cos\beta+\sin\alpha\odot\sin\beta,
  \label{eq:cos-sub}\\
  \sin(\beta-\alpha)&=\sin\beta\odot\cos\alpha-\cos\beta\odot\sin\alpha,
  \label{eq:sin-sub}
\end{align}
which are
\lmod{Tensor.CosSub.eq.AddMulS} and
\lmod{Tensor.SinSub.eq.SubMulSSin_Cos}.

The RoFormer matrix satisfies the same law by conjugation.

\begin{theorem}
\label{thm:dotTp}
\lmod{Tensor.DotT.eq.RotaryMatrix'Sub}~\cite{lemma_rope_su_sub}.
For all \(\alpha,\beta\in\R^d\),
\begin{equation}
  \rot'(\alpha)^\top @ \rot'(\beta) = \rot'(\beta-\alpha).
\end{equation}
\end{theorem}

The proof writes \(\rot'(\alpha)=\boldsymbol{P}^\top\rot(\alpha)\boldsymbol{P}\) (Theorem~\ref{thm:conj}), cancels \(\boldsymbol{P}^\top@\boldsymbol{P}=\eye\) and \(\boldsymbol{P}@\boldsymbol{P}^\top=\eye\), and applies Theorem~\ref{thm:dotT}.

\begin{theorem}[Additive group law]
\label{thm:add}
\lmod{Tensor.DotRotaryMatrixS.eq.RotaryMatrixAdd}~\cite{lemma_rope_add}.
For all \(\alpha,\beta\in\R^d\),
\begin{equation}
  \rot(\alpha)\,@\,\rot(\beta) = \rot(\alpha+\beta).
\end{equation}
\end{theorem}

\begin{proof}
Rewrite \(\alpha+\beta=\beta-(-\alpha)\)
(\texttt{add\_comm}, \lmod{Int.Add.eq.Sub_Neg})
and \(\alpha=-(-\alpha)\).
Lemma~\ref{lem:neg-T} gives \(\rot(\alpha)=\rot(-\alpha)^\top\).
Theorem~\ref{thm:dotT} then yields
\(\rot(-\alpha)^\top @ \rot(\beta)=\rot(\beta-(-\alpha))\).
\end{proof}

Theorem~\ref{thm:add} is the split-half, free-angle form of the additivity already stated by Liu et al.~\cite{liu2025rethinking} and Zhang et al.~\cite{zhang2025grape}.
It holds here for a full split-half matrix and an arbitrary pair of angle vectors: there is no token index and no geometric schedule.
The special case \(\beta=\alpha\) recovers Theorem~\ref{thm:orth} via \(\rot(0)=\eye\).

Composing Theorem~\ref{thm:dotT} with linearity of \(\theta\) in the index produces the token-indexed form.

\begin{proposition}
\label{prop:eqdott}
\lmod{Tensor.DotT.eq.RotaryMatrixSub.of.Eq_Stack_Mul.Ge}.
Suppose \(\theta\in\R^{n\times d}\) satisfies \(\theta_i=i\,\tau\) for a fixed \(\tau\in\R^d\), and let \(k,t:\mathrm{Fin}\,n\) with \(k\ge t\).
Then
\begin{equation}
  \rot(\theta_t)^\top @ \rot(\theta_k) = \rot(\theta_{k-t}).
\end{equation}
The same arithmetic on the RoFormer layout is
\lmod{Tensor.DotT.eq.RotaryMatrix'Get_Sub.of.Eq_Stack_Mul.Ge}:
\(\rot'(\theta_t)^\top @ \rot'(\theta_k) = \rot'(\theta_{k-t})\).
\end{proposition}

The only arithmetic beyond Theorem~\ref{thm:dotT} is
\(\theta_k-\theta_t=\theta_{k-t}\)
(\lmod{Tensor.SubGetS.eq.Get_Sub.of.Eq_Stack_Mul.Ge}),
which is \((k-t)\tau\) once \(k-t\) is read as a value rather than a wrap
(\lmod{Nat.CoeSub.eq.SubCoeS.of.Ge}).

\section{Relative position in attention}
\label{sec:rel}

\subsection{Pairing of two rotated vectors}
\label{sec:pair}

RoFormer's relative-position law is the inner-product identity
\begin{equation}
  \label{eq:roformer-rel-again}
  \bigl(\rot_{\Theta,m}\,q\bigr)^\top\bigl(\rot_{\Theta,n}\,k\bigr)
  = q^\top \rot_{\Theta,n-m}\,k.
\end{equation}
The library proves this for free angle vectors, with no frequency hypothesis.
On two \(1\)-tensors the library writes the inner product as \(@\), so~\eqref{eq:roformer-rel-again} is definitionally
\begin{equation}
  \label{eq:pair-lean}
  \bigl(\rot(\alpha)\,@\,q\bigr)\,@\,\bigl(\rot(\beta)\,@\,k\bigr)
  \;=\;
  q\,@\,\bigl(\rot(\beta-\alpha)\,@\,k\bigr).
\end{equation}
The correspondence is \(\alpha\leftrightarrow\Theta_m\), \(\beta\leftrightarrow\Theta_n\), hence \(\beta-\alpha\leftrightarrow n-m\) in angle space.
The Lean~4 statement, with a complete proof and no \texttt{sorry}, is
\begin{center}
\lmod{Tensor.DotDotSRotaryMatrix.eq.Dot_DotRotaryMatrixSub}.
\end{center}

\begin{theorem}[RoFormer pairing]
\label{thm:pair}
\lmod{Tensor.DotDotSRotaryMatrix.eq.Dot_DotRotaryMatrixSub}~\cite{lemma_rope_pair}.
For \(\alpha,\beta\in\R^d\) and \(q,k\in\R^{d+d}\),
equation~\eqref{eq:pair-lean} holds.
\end{theorem}

\begin{proof}
The Lean proof is the following chain.
First, a vector--matrix product may be moved across the inner product
(\lmod{Tensor.Dot_T.eq.Dot}):
\begin{equation}
  \label{eq:pair-step1}
  (\rot(\alpha)\,@\,q)\,@\,(\rot(\beta)\,@\,k)
  \;=\;
  \bigl(q\,@\,\rot(\alpha)^\top\bigr)\,@\,(\rot(\beta)\,@\,k).
\end{equation}
Associativity of \(@\) in the vector--matrix--vector case
(\lmod{Tensor.DotDot.eq.Dot_Dot}, instance \texttt{vmv}) then gives
\begin{equation}
  \label{eq:pair-step2}
  \bigl(q\,@\,\rot(\alpha)^\top\bigr)\,@\,(\rot(\beta)\,@\,k)
  \;=\;
  q\,@\,\bigl(\rot(\alpha)^\top @ (\rot(\beta)\,@\,k)\bigr).
\end{equation}
A second associativity
(\lmod{Tensor.DotDot.eq.Dot_Dot}, instance \texttt{mmv}) pulls the two rotary matrices together:
\begin{equation}
  \label{eq:pair-step3}
  \rot(\alpha)^\top @ (\rot(\beta)\,@\,k)
  \;=\;
  \bigl(\rot(\alpha)^\top @ \rot(\beta)\bigr)\,@\,k.
\end{equation}
Theorem~\ref{thm:dotT} replaces the product of matrices by a single rotary matrix of the angle difference,
\begin{equation}
  \label{eq:pair-step4}
  \rot(\alpha)^\top @ \rot(\beta) = \rot(\beta-\alpha).
\end{equation}
Substituting~\eqref{eq:pair-step4} into~\eqref{eq:pair-step3} and then into~\eqref{eq:pair-step2} yields~\eqref{eq:pair-lean}.
\end{proof}

No token index and no schedule on \(\Theta\) appear in Theorem~\ref{thm:pair}.
The classical notation \(\rot_{\Theta,n-m}\) is recovered in the next subsection, once \(\theta_i=i\,\tau\) identifies \(\beta-\alpha\) with a row of \(\theta\).

\subsection{Linear frequencies}

Relative \emph{indexing}---replacing \(\theta_k-\theta_i\) by a row of \(\theta\) itself---needs a relation among the rows of \(\theta\).

\begin{definition}[Linear frequency hypothesis]
\label{def:lin}
A tensor \(\theta\in\R^{n\times d}\) is linear in the token index when there exists \(\tau\in\R^d\) such that
\begin{equation}
  \theta = \bigl[\,i < n\,\bigr]\;(\tau\cdot i)
  \qquad\text{i.e.}\qquad
  \theta_i = i\,\tau\in\R^d.
\end{equation}
In Lean this is the hypothesis
\texttt{h\(\theta\) : \(\theta\) = [i < n] (\(\tau\) * (i : \(\mathbb{R}\)))}.
\end{definition}

Standard RoPE is the instance
\(\tau_p=\lambda/b^{p/d}\).
The lemmas below never mention \(b\) or \(\lambda\).

\begin{lemma}
\label{lem:ite}
\lmod{Tensor.RotaryMatrixSubGetS.eq.Ite_RotaryMatrix_T.of.Eq_Stack_Mul}.
Under Definition~\ref{def:lin}, for all \(i,j:\mathrm{Fin}\,n\),
\begin{equation}
  \rot(\theta_j-\theta_i)
  \;=\;
  \begin{cases}
    \rot(\theta_{j-i}) & \text{if }j\ge i,\\[0.25em]
    \rot(\theta_{i-j})^\top & \text{if }j<i.
  \end{cases}
\end{equation}
\end{lemma}

The non-negative branch is \(\theta_j-\theta_i=\theta_{j-i}\).
The negative branch is \(\theta_j-\theta_i=-(\theta_i-\theta_j)=-\theta_{i-j}\) and Lemma~\ref{lem:neg-T}.
The library proof uses \texttt{erw [if\_pos]} / \texttt{erw [if\_neg]}: a plain \texttt{rw} does not fire on the \texttt{GetElem} instance of \(\le\) on \(\mathrm{Fin}\,n\).

\begin{remark}
Necessity, which we do not formalize as a separate lemma, is slightly weaker: \(\rot\) depends on an angle only through \((\cos,\sin)\), so it is enough that
\(\theta_j-\theta_i\equiv\theta_{j-i}\pmod{2\pi}\)
componentwise on \(j\ge i\).
Definition~\ref{def:lin} is the clean generator of that congruence.
We kept equality of tensors, not congruence, to avoid a side development of \(2\pi\mathbb{Z}\).
\end{remark}

\subsection{Softmax attention}

Let \(Q,K,V\in\R^{n\times(d+d)}\) and write
\(R(i):=\rot(\theta_i)\).
The RoPE attention map is
\begin{equation}
  \label{eq:rope-attn}
  \softmax\!\left(\frac{Q^R (K^R)^\top}{\sqrt{d+d}}\right) @ V,
  \qquad
  Q^R_i = R(i)\,@\,Q_i,\quad
  K^R_i = R(i)\,@\,K_i.
\end{equation}

\begin{theorem}[Relative softmax]
\label{thm:soft}
\lmod{Tensor.DotSoftmaxDivDot_Stack_TDot.eq.Stack_Div_SumExp.of.Eq_Stack_Mul}~\cite{lemma_rope_rel}.
Under Definition~\ref{def:lin}, the map~\eqref{eq:rope-attn} equals
\begin{equation}
  \label{eq:soft}
  \Biggl[
    i < n,\;
    j < d+d
  \Biggr]
  \frac{
    \displaystyle\sum_{k}
      V_{k,j}\,
      \exp\Bigl(
        Q_i @ \bigl(\Rel(i,k)\,@\,K_k\bigr)
        \big/\sqrt{d+d}
      \Bigr)
  }{
    \displaystyle
      \sum\exp\Bigl(
        \bigl(R(i)\,@\,Q_i\bigr)
        @
        (K^R)^\top
        \big/\sqrt{d+d}
      \Bigr)
  },
\end{equation}
where the relative matrix is the one from Lemma~\ref{lem:ite},
\begin{equation}
  \label{eq:rel}
  \Rel(i,k)
  \;=\;
  \begin{cases}
    R(k-i) & \text{if }k\ge i,\\
    R(i-k)^\top & \text{if }k<i,
  \end{cases}
\end{equation}
and the denominator is the sum of the exponential of the full logit row of query \(i\) (shape \([n]\)), left in rotated form.
\end{theorem}

The proof is a stack-wise softmax expansion
(\lmod{Tensor.DotSoftmaxDivDot_T.eq.Stack_Div_SumExp})
followed, in each numerator summand, by Theorem~\ref{thm:pair} on the pair \((i,k)\) and Lemma~\ref{lem:ite} on the resulting angle difference.
No extra hypothesis on \(Q,K,V\) is required.

Two bookkeeping points match the Lean statement exactly.
First, \(\Rel\) is bound in terms of \(R\), not by a second mention of \texttt{rotaryMatrix}, so the relative matrix is definitionally a rotary matrix or its transpose.
Second, the partition function is not rewritten with \(\Rel\): it remains
\(\sum\exp\bigl((R(i)@Q_i)@(K^R)^\top/\sqrt{d+d}\bigr)\).
The two forms are equal by the same pairing, but the certified identity only rewrites the weights that multiply \(V\).

Equation~\eqref{eq:soft} is the sense in which ``RoPE is a relative position embedding'': after the rotations are applied, every unnormalized weight is an ordinary inner product of the unrotated query against a key that has been rotated by a matrix depending only on \(k-i\).

\section{Library structure}
\label{sec:artifact}

Interactive visualizations of the public statements:
\begin{itemize}[noitemsep]
  \item Relative softmax:
    \lmod{Tensor.DotSoftmaxDivDot_Stack_TDot.eq.Stack_Div_SumExp.of.Eq_Stack_Mul}
  \item RoFormer pairing~\eqref{eq:roformer-rel}:
    \lmod{Tensor.DotDotSRotaryMatrix.eq.Dot_DotRotaryMatrixSub}
  \item Conjugation \(R'=P^\top RP\):
    \lmod{Tensor.RotaryMatrix'.eq.DotDot_RotaryMatrix}
  \item Hugging Face matrix:
    \lmod{Tensor.RotaryMatrix.eq.AppendHstackSMulSEye}
  \item RoFormer matrix:
    \lmod{Tensor.RotaryMatrix'.eq.Stack_Ite_IteS}
  \item Realization:
    \lmod{Tensor.DotRotaryMatrix.eq.AddMulS}
  \item Additive group law:
    \lmod{Tensor.DotRotaryMatrixS.eq.RotaryMatrixAdd}
  \item Orthogonality:
    \lmod{Tensor.DotT_RotaryMatrix.eq.Eye}
\end{itemize}
Source: \href{https://github.com/math-proof/lemma}{github.com/math-proof/lemma}~\cite{lemma_repo}.
Every module in Table~\ref{tab:mods} builds with \texttt{lake build} and contains no \texttt{sorry}.

\begin{table}[t]
\centering
\caption{Lean modules for one-dimensional RoPE.
The public name of a lemma is the module path.}
\label{tab:mods}
\small
\begin{tabular}{p{0.62\textwidth}p{0.30\textwidth}}
\hline
Module & Content \\
\hline
\texttt{RotaryMatrix.eq.AppendHstackSMulSEye} & Definition~\ref{def:rot} (Hugging Face) \\
\texttt{RotaryMatrix'.eq.Stack\_Ite\_IteS} & Definition~\ref{def:rotp} (RoFormer) \\
\texttt{Interleave.eq.AppendStackS\_Delta} & Definition~\ref{def:P} \\
\texttt{RotaryMatrix'.eq.DotDot\_RotaryMatrix} & Theorem~\ref{thm:conj} \\
\texttt{GetRotaryMatrix'.eq.GetRotaryMatrix} & Theorem~\ref{thm:entry} \\
\texttt{DotRotaryMatrix.eq.AddMulS} & Theorem~\ref{thm:real} \\
\texttt{RotaryMatrixNeg.eq.TRotaryMatrix} & Lemma~\ref{lem:neg-T} \\
\texttt{RotaryMatrix0.eq.Eye} & Lemma~\ref{lem:R0} \\
\texttt{RotaryMatrix'0.eq.Eye} & \(\rot'(0)=\eye\) \\
\texttt{DotT\_RotaryMatrix.eq.Eye} & Theorem~\ref{thm:orth} \\
\texttt{DotRotaryMatrixNeg.eq.Eye} & \(\rot(-\alpha)\rot(\alpha)=\eye\) \\
\texttt{DotT.eq.RotaryMatrixSub} & Theorem~\ref{thm:dotT} \\
\texttt{DotT.eq.RotaryMatrix'Sub} & Theorem~\ref{thm:dotTp} \\
\texttt{DotRotaryMatrixS.eq.RotaryMatrixAdd} & Theorem~\ref{thm:add} \\
\texttt{SubGetS.eq.Get\_Sub.of.Eq\_Stack\_Mul.Ge} & \(\theta_k-\theta_t=\theta_{k-t}\) \\
\texttt{DotT.eq.RotaryMatrixSub.of.Eq\_Stack\_Mul.Ge} & Proposition~\ref{prop:eqdott} \\
{\footnotesize\lmod{Tensor.DotT.eq.RotaryMatrix'Get_Sub.of.Eq_Stack_Mul.Ge}} & \(\rot'\) form of Proposition~\ref{prop:eqdott} \\
{\footnotesize\lmod{Tensor.RotaryMatrixSubGetS.eq.Ite_RotaryMatrix_T.of.Eq_Stack_Mul}} & Lemma~\ref{lem:ite} \\
\texttt{DotDotSRotaryMatrix.eq.Dot\_DotRotaryMatrixSub} & Theorem~\ref{thm:pair} \\
{\footnotesize\lmod{Tensor.DotSoftmaxDivDot_Stack_TDot.eq.Stack_Div_SumExp.of.Eq_Stack_Mul}} & Theorem~\ref{thm:soft} \\
\hline
\end{tabular}
\end{table}

The cosine and sine facts used above are proved first on \(\R\), then on \texttt{List.Vector}, then on tensors by transport along \texttt{.data}.
That stratification is deliberate: the rotary proofs never unfold the tensor implementation past the corresponding vector lemma.

Figure~\ref{fig:deps} lists the import graph of the headline modules.

\begin{figure}[ht]
\footnotesize
\begin{verbatim}
DotSoftmaxDivDot_Stack_TDot.eq.Stack_Div_SumExp.of.Eq_Stack_Mul
|-- DotSoftmaxDivDot_T.eq.Stack_Div_SumExp
|-- DotDotSRotaryMatrix.eq.Dot_DotRotaryMatrixSub
|   `-- DotT.eq.RotaryMatrixSub
`-- RotaryMatrixSubGetS.eq.Ite_RotaryMatrix_T.of.Eq_Stack_Mul
    |-- SubGetS.eq.Get_Sub.of.Eq_Stack_Mul.Ge
    `-- RotaryMatrixNeg.eq.TRotaryMatrix

DotRotaryMatrixS.eq.RotaryMatrixAdd
|-- DotT.eq.RotaryMatrixSub
`-- RotaryMatrixNeg.eq.TRotaryMatrix

DotT_RotaryMatrix.eq.Eye
|-- DotT.eq.RotaryMatrixSub
`-- RotaryMatrix0.eq.Eye

DotRotaryMatrix.eq.AddMulS
`-- DotRotaryMatrix.eq.AddMulSAppend

RotaryMatrix'.eq.DotDot_RotaryMatrix
|-- GetRotaryMatrix'.eq.Ite_IteS
`-- GetRotaryMatrix.eq.Mul{Cos,NegSin,Sin}_Delta (four quadrants)

DotT.eq.RotaryMatrix'Sub
|-- RotaryMatrix'.eq.DotDot_RotaryMatrix
|-- DotT.eq.RotaryMatrixSub
|-- DotTInterleave.eq.Eye
`-- Dot_TInterleave.eq.Eye
\end{verbatim}
\caption{Import graph for the headline RoPE lemmas.
Prefixes \texttt{Tensor.} are omitted.}
\label{fig:deps}
\end{figure}

\section{Discussion and limitations}
\label{sec:discussion}

\paragraph{What is formalized, and what is not.}
We certify equalities of tensors, not approximation quality, length extrapolation, or the interaction of RoPE with trained weights.
Theorem~\ref{thm:soft} is an identity for every \(Q,K,V\) and every linear \(\theta\); it does not say that a trained model uses the relative matrix in a particular way at inference time, nor that attention decays with distance~\cite{barbero2025round}.

\paragraph{Layout.}
We formalize both pairings.
\texttt{rotaryMatrix} is the Hugging Face split-half matrix~\eqref{eq:rot-hf};
\texttt{rotaryMatrix'} is Su et al.'s interleaved matrix~\eqref{eq:rot-su}.
Theorem~\ref{thm:conj} is the permutation lemma that earlier drafts left informal.
Kernel realization (Theorem~\ref{thm:real}) and the inner-product pairing (Theorem~\ref{thm:pair}) are stated for the Hugging Face layout, which is the map \texttt{rotate\_half} implements.
The matching matrix identity on the RoFormer layout is Theorem~\ref{thm:dotTp}.

\paragraph{Frequencies.}
Definition~\ref{def:lin} is sufficient for Lemma~\ref{lem:ite} and strictly weaker than the classical geometric schedule.
It is not necessary: any \(\theta\) with
\(\theta_k-\theta_t\equiv\theta_{k-t}\pmod{2\pi}\)
gives the same \(\rot\).

\paragraph{Two-dimensional and multimodal RoPE.}
Axial or multi-axis rotary embeddings reuse Definition~\ref{def:rot} on several index streams.
The additive and orthogonal theorems apply unchanged to each stream; a relative-index lemma needs a linearity hypothesis per axis.
We leave that packaging to later work.

\paragraph{Typing of the scalar inner product.}
A one-dimensional \texttt{@} elaborates as a tensor whose shape is
\(\mathrm{matmul\_shape}\,[d]\,[d]\),
which is not definitionally \(\Tensor\,\R\,[]\).
The Lean statements therefore wrap those scalars in
\texttt{id ($\alpha$ := Tensor \(\R\) [])}.
This is a library elaboration fact, not a mathematical hypothesis.

\section{Conclusion}
\label{sec:conclusion}

The algebraic skeleton of one-dimensional RoPE is a handful of identities about a block matrix of plane rotations, in either the Hugging Face split-half layout or Su et al.'s interleaved layout.
We have stated both matrices in Lean~4, proved they are conjugate by a fixed even/odd gather, and, on the Hugging Face matrix that implementations multiply, proved the kernel evaluation formula, the orthogonal and additive group laws for free angle vectors, the RoFormer pairing~\eqref{eq:roformer-rel}, and the relative-offset form of softmax attention under a linear frequency hypothesis.
The modules named in the introduction are the public interface of that development~\cite{lemma_rope_rel,lemma_rope_pair,lemma_rope_conj,lemma_rope_real,lemma_rope_add,lemma_rope_orth}.

\section*{Acknowledgements}

The formalization uses Lean~4 and mathlib~\cite{moura2021lean4,mathlib2020}.
Interactive statements of the modules in Table~\ref{tab:mods} are available at \href{http://www.lemma.cn/}{lemma.cn}.

\bibliographystyle{plain}
\bibliography{refs}

\end{document}
